% ============================================================================
% figures/sequence_chat.tex
% Diagramme de séquence d'un échange Q&A en streaming.
% Lifelines positionnées manuellement.
% ============================================================================
\begin{figure}[H]
\centering
\begin{tikzpicture}[
  font=\footnotesize\sffamily,
  every node/.style={align=center, inner sep=3pt},
  actor/.style={draw=black!40, fill=black!5, rounded corners=2pt, minimum width=2.2cm, minimum height=0.7cm},
  lifeline/.style={dashed, gray!60},
  msg/.style={-Stealth, thick, blue!40!black},
  ret/.style={-Stealth, thick, dashed, green!50!black},
]

% Positions X des acteurs
\def\ux{0}
\def\nx{3}
\def\fx{6}
\def\rx{9}
\def\cx{12}
\def\ox{14.8}

% Acteurs (tête)
\node[actor] (U) at (\ux, 0) {\textbf{Utilisateur}};
\node[actor] (N) at (\nx, 0) {\textbf{Next.js}};
\node[actor] (F) at (\fx, 0) {\textbf{FastAPI}};
\node[actor] (R) at (\rx, 0) {\textbf{Routeur}};
\node[actor] (C) at (\cx, 0) {\textbf{ChromaDB}};
\node[actor] (O) at (\ox, 0) {\textbf{OpenAI}};

% Lifelines
\foreach \x/\name in {\ux/U, \nx/N, \fx/F, \rx/R, \cx/C, \ox/O}
  \draw[lifeline] (\x, -0.35) -- (\x, -8.5);

% Pied d'acteurs
\foreach \x/\name in {\ux/U, \nx/N, \fx/F, \rx/R, \cx/C, \ox/O}
  \node[actor] at (\x, -8.7) {\textbf{\name}};

% Échanges (de haut en bas)
\draw[msg] (\ux, -0.9) -- node[above, font=\scriptsize] {question tapée} (\nx, -0.9);
\draw[msg] (\nx, -1.6) -- node[above, font=\scriptsize] {POST /api/chat/stream} (\fx, -1.6);
\draw[msg] (\fx, -2.3) -- node[above, font=\scriptsize] {détection intention} (\rx, -2.3);
\draw[msg] (\rx, -3.0) -- node[above, font=\scriptsize] {recherche top-5} (\cx, -3.0);
\draw[ret] (\cx, -3.6) -- node[above, font=\scriptsize] {5 chunks + métadonnées} (\rx, -3.6);
\draw[msg] (\rx, -4.3) -- node[above, font=\scriptsize] {prompt + contexte} (\ox, -4.3);

% Boucle streaming
\draw[dotted, thick, gray] (\rx-0.5, -4.8) -- (\ox+0.5, -4.8);
\node[anchor=west, font=\scriptsize\itshape] at (\ox+0.6, -4.8) {loop};

\draw[ret] (\ox, -5.5) -- node[above, font=\scriptsize] {token} (\rx, -5.5);
\draw[ret] (\rx, -6.2) -- node[above, font=\scriptsize] {SSE: event token} (\fx, -6.2);
\draw[ret] (\fx, -6.8) -- node[above, font=\scriptsize] {chunk HTTP} (\nx, -6.8);
\draw[ret] (\nx, -7.4) -- node[above, font=\scriptsize] {affichage incrémental} (\ux, -7.4);

% Fin
\draw[ret] (\ox, -8.0) -- node[above, font=\scriptsize] {done} (\rx, -8.0);

\end{tikzpicture}
\caption{Diagramme de séquence d'un échange Q\&A en \textit{streaming}. Après la phase de \textit{retrieval} dans ChromaDB, le routeur relaie la requête au \ac{LLM} et retransmet chaque \textit{token} reçu à l'interface via un flux \ac{SSE}.}
\label{fig:sequence_chat}
\end{figure}
